% TOP_PROBLEM: {"id": 2540, "title": "Kourovka Notebook Problem 21.31", "category_id": 17, "category": {"id": 17, "name": "group_theory", "display_name": "Group Theory", "description": "Problems about groups, group actions, representations, and related algebraic structures.", "slug": "group-theory", "order_index": 17, "created_at": "2026-07-31T15:26:25.671Z"}, "status": "open", "problem_number": "KOU-21.31", "set_id": 10}
% FIRST_POSTED: 2026-10-01
% ENTRY_KIND: statement_only
% UnsolvedMath ID 2540; source catalogue code KOU-21.31
% !TeX program = lualatex
% Definition and sources only; this file is not an LLM solution attempt.
\documentclass[11pt,a4paper]{article}
\usepackage[margin=25mm]{geometry}
\usepackage{fontspec}
\setmainfont{lmroman10-regular.otf}[BoldFont=lmroman10-bold.otf,
 ItalicFont=lmroman10-italic.otf,BoldItalicFont=lmroman10-bolditalic.otf]
\usepackage{amsmath,amssymb,mathtools}
\usepackage{unicode-math}
\setmathfont{latinmodern-math.otf}
\AtBeginDocument{\let\setminus\smallsetminus}
\usepackage{xurl,hyperref}
\hypersetup{colorlinks=true,urlcolor=blue}
\setcounter{secnumdepth}{0}
\setlength{\parindent}{0pt}
\setlength{\parskip}{5pt}
\setlength{\emergencystretch}{3em}
\begin{document}
\section*{Kourovka Notebook Problem 21.31}\label{2540}
{\small UnsolvedMath ID 2540 · KOU-21.31 · Group Theory · Kourovka Notebook - New Problems, Issue 21}
\subsection{Definitions and mathematical statement}
Let $N$ be a finite group and let $\operatorname{Aut}(N)$ be its group
of automorphisms. Its holomorph is the semidirect product
\[
 \operatorname{Hol}(N)=N\rtimes\operatorname{Aut}(N),\qquad
 (n,\alpha)(m,\beta)=(n\alpha(m),\alpha\beta).
\]
It acts on the underlying set of $N$ by $(n,\alpha)\cdot x=n\alpha(x)$.
This action is faithful. A subgroup $R\le\operatorname{Hol}(N)$ is
regular if for every $x,y\in N$ exactly one element of $R$ sends $x$
to $y$. Equivalently, the action is transitive and every point stabilizer
is trivial; hence $|R|=|N|$.

A group $K$ is soluble if its derived series $K^{(0)}=K$,
$K^{(i+1)}=[K^{(i)},K^{(i)}]$ reaches $1$, where brackets mean the
subgroup generated by commutators.
The conjecture states
\[
 N\text{ finite soluble},\quad R\le\operatorname{Hol}(N)\text{ regular}
 \quad\Longrightarrow\quad R\text{ soluble}.
\]
The subgroup $R$ is not assumed normal in the holomorph, and need not
be its evident translation subgroup $N\times\{\mathrm{id}\}$.
Its multiplication is the multiplication inherited from the holomorph,
not a multiplication transported from $N$ by an arbitrary set bijection.

Regularity specifies a particularly restrictive embedding of a group of
order $|N|$ into the permutation group of $N$. The question does not
claim that every group of that order is soluble. Both the algebraic
solubility assumption and the regularity of this specific action must
hold in each instance.
\subsection{Short English statement}
Must every regular subgroup of the holomorph of a finite soluble group itself be soluble?
\subsection{Sources}
\begin{itemize}
\item[S1] ULAM.AI and the UnsolvedMath Contributors, supplied problems.json, record 2540 (KOU-21.31), Kourovka Notebook - New Problems, Issue 21. Catalogue identity and source transcription; CC BY 4.0.
\url{https://huggingface.co/datasets/ulamai/UnsolvedMath}.
\item[S2] The Kourovka Notebook, 21st edition, new problems, Problem 21.31. Expanded formulation of the supplied catalogue statement.
\url{https://arxiv.org/pdf/1401.0300v47}.
\end{itemize}
\end{document}
